\newpage
\chapter{RESULTADOS ESPERADOS}

Al concluir el proyecto, se espera que la plataforma web desarrollada genere un impacto
medible tanto en la eficiencia de la gestión administrativa del TUPA como en la experiencia
de los distintos actores de la comunidad universitaria. Los resultados se clasifican en
beneficios cuantitativos y cualitativos, y se alinean directamente con los ocho objetivos

\section{Resultados cuantitativos}
Los siguientes indicadores de resultado proyectan el impacto esperado de la solución una vez
desplegada en producción y adoptada por la comunidad universitaria:

{\small
\begin{tabular}{|p{3.9cm}|c|c|p{2.6cm}|}
\hline
\textbf{Resultado esperado} & \textbf{Línea base (est.)} & \textbf{Meta} & \textbf{OE relacionado} \\
\hline
Reducción del tiempo promedio de atención por trámite & 5–7 días hábiles & $\leq$ 3 días hábiles & OE-03, OE-04 \\
\hline
Reducción de consultas presenciales innecesarias en ventanilla & 100\% presencial & $\leq$ 30\% presencial & OE-01, OE-02 \\
\hline
Incremento en la tasa de trámites resueltos en primer intento & $\sim$40\% & $\geq$ 75\% & OE-03, OE-05 \\
\hline
Porcentaje de módulos con panel de control operativo & 0\% (no existe) & 100\% & OE-06 \\
\hline
Disponibilidad del sistema en producción & N/A & $\geq$ 99\% mensual & OE-08 \\
\hline
Tiempo de respuesta del sistema (p95) & N/A & $\leq$ 2 segundos & OE-01, OE-08 \\
\hline
Generación automática de reportes administrativos & Manual (0 automatiz.) & $\geq$ 5 tipos de reporte & OE-07 \\
\hline
\end{tabular}
}

\textit{Nota: Los valores de línea base son estimaciones obtenidas del análisis de la
problemática descrita. Serán validados con las oficinas administrativas
durante la fase de levantamiento de requerimientos.}

\section{Resultados cualitativos}

\begin{itemize}
    \item \textbf{Mejora de la experiencia de usuario:} La plataforma ofrecerá una interfaz
    moderna, intuitiva y responsiva, reduciendo la curva de aprendizaje para estudiantes y
    personal administrativo sin conocimientos técnicos avanzados (alineado con OE-01).

    \item \textbf{Mayor transparencia institucional:} Los estudiantes podrán consultar en tiempo
    real el estado de sus trámites, los requisitos exigidos y los costos vigentes, eliminando
    la asimetría de información característica del sistema actual (OE-02, OE-04).

    \item \textbf{Centralización y consistencia de la información:} La solución integrará los
    procedimientos de distintas oficinas en un único punto de acceso, evitando duplicidades
    y versiones desactualizadas del TUPA (OE-02, OE-06).

    \item \textbf{Reducción de la sobrecarga administrativa:} La automatización de notificaciones
    y validaciones documentales liberará tiempo del personal de ventanilla para tareas de
    mayor valor agregado (OE-04, OE-05).

    \item \textbf{Cultura de toma de decisiones basada en datos:} La generación de reportes y
    estadísticas de trámites procesados facilitará la identificación de cuellos de botella y
    la mejora continua de los procesos administrativos (OE-07).

    \item \textbf{Sostenibilidad y escalabilidad técnica:} El despliegue en infraestructura en la
    nube y el uso de estándares de desarrollo garantizarán la mantenibilidad y posibilidad de
    extensión futura de la plataforma (OE-08).
\end{itemize}

\section{Impacto esperado}

{\small
\begin{tabular}{|p{2.8cm}|p{5.0cm}|p{4.8cm}|}
\hline
\textbf{Actor} & \textbf{Beneficio directo} & \textbf{Indicador de impacto} \\
\hline
Estudiantes & Acceso 24/7 a información del TUPA, registro en línea de solicitudes y seguimiento del estado de trámites & Reducción de tiempos de espera; disminución de traslados presenciales \\
\hline
Personal administrativo & Validación digital de documentos, panel de gestión por oficina, notificaciones automáticas de acciones pendientes & Menor carga manual; reducción de errores en el proceso \\
\hline
Jefes de oficina & Acceso a reportes estadísticos de trámites, monitoreo del desempeño del área y gestión de usuarios del sistema & Mejor visibilidad para toma de decisiones gerenciales \\
\hline
UNSAAC (institución) & Modernización de la imagen institucional, reducción de costos operativos y cumplimiento de estándares de gobierno digital & Mejora en indicadores de gestión universitaria \\
\hline
\end{tabular}
}

\section{Alineación con los Objetivos Específicos del Proyecto}

La siguiente matriz evidencia la correspondencia directa entre los resultados esperados y los
ocho objetivos específicos del proyecto, garantizando la trazabilidad requerida por el indicador
de desempeño \textbf{AG-C12.01}:

{\small
\begin{tabular}{|p{1.0cm}|p{4.8cm}|p{6.4cm}|}
\hline
\textbf{OE} & \textbf{Objetivo Específico} & \textbf{Resultado / Entregable asociado} \\
\hline
OE-01 & Mejorar la interfaz y experiencia de usuario & Rediseño UX/UI; puntuación de usabilidad $\geq$ 4.0/5.0 \\
\hline
OE-02 & Facilitar la consulta de procedimientos, requisitos y costos & Módulo de búsqueda y consulta pública del TUPA \\
\hline
OE-03 & Optimizar el registro y seguimiento de trámites & Sistema de gestión de solicitudes con trazabilidad por estado \\
\hline
OE-04 & Implementar notificaciones sobre el estado de los procedimientos & Motor de notificaciones (correo electrónico / en plataforma) \\
\hline
OE-05 & Mejorar la gestión y validación de documentos digitales & Módulo de carga, validación y aprobación de documentos \\
\hline
OE-06 & Panel de control para oficinas administrativas & Dashboard por rol con métricas en tiempo real \\
\hline
OE-07 & Generación de reportes y estadísticas & Motor de reportes exportable (PDF / Excel) \\
\hline
OE-08 & Garantizar accesibilidad y disponibilidad & Despliegue en la nube con SLA $\geq$ 99\%; diseño responsivo \\
\hline
\end{tabular}
}


